\documentclass[11pt,a4paper]{article}

\usepackage[utf8]{inputenc}

\usepackage[T1]{fontenc}

\usepackage[polish]{babel}

\usepackage{lmodern}

\usepackage{geometry}

\geometry{margin=27mm}

\usepackage[hidelinks]{hyperref}

\title{Odbiór zachowania humanoida na nagraniu wideo}

\author{Mateusz Pomianek}

\date{Szkic naukowy, wersja 1.0.0 | 28 listopada 2026}

\begin{document}

\maketitle

\begin{abstract}

Proponowana miara główna to różnica prawdopodobieństw prawidłowego odgadnięcia celu ruchu.

\end{abstract}

\section{Wprowadzenie}

Ocena sprawczości wiąże się z przypisywaniem zdolności do zamierzonego działania (Gray et al., 2007). W animacji ruch przygotowawczy służy wskazaniu następującego po nim działania (Lasseter, 1987).

\section{Cel i pytania badawcze}

Hipotezy robocze wymagają zatwierdzenia przed rozpoczęciem właściwej weryfikacji.

\paragraph{H1.} Klipy z ruchem przygotowawczym zgodnym z późniejszym celem zwiększą prawdopodobieństwo poprawnego odgadnięcia celu względem klipów bez informacji kierunkowej, gdy nagranie zostanie zatrzymane przed ruchem właściwym.

\section{Materiał i metody}

Pierwszorzędowy wskaźnik: różnica prawdopodobieństw prawidłowego odgadnięcia celu ruchu.

Warianty tego samego ruchu zachowają kadr, czas i prędkość właściwego odcinka. Wynik główny to różnica prawdopodobieństw poprawnej odpowiedzi z modelu logistycznego krzyżującego osobę i klip. Ocena subiektywnej przewidywalności i intencjonalności będzie wynikiem dodatkowym.

Przed badaniem z udziałem ludzi należy zatwierdzić wymagania etyczne, zasady zgody i ochronę danych.

\section{Plan analizy}

Zamrożenie estymandu, reguł oceny jakości i obsługi braków nastąpi przed dostępem do danych wynikowych. Raport obejmie oszacowanie kontrastu głównego z przedziałem niepewności i pełne mianowniki. Jeśli przewidziano hipotezy drugorzędowe, przed analizą należy ustalić sposób kontroli wielokrotności.

\section{Wyniki}

W niniejszym pliku nie umieszczono wyników empirycznych, statystyk ani nieprzeprowadzonych analiz. Sekcja wymaga uzupełnienia po uzyskaniu i weryfikacji materiału dowodowego.

\section{Dyskusja i ograniczenia}

Ograniczenia obejmują możliwe błędy pomiaru, dobór próby, alternatywne mechanizmy i ograniczoną możliwość uogólniania. Korelacja nie dowodzi przyczynowości; poprawność techniczna nie jest dowodem znaczenia klinicznego.

\section{Wnioski}

Hipotezy i sposób pomiaru stanowią podstawę oceny wykonalności oraz przygotowania właściwego protokołu.

\section*{Bibliografia}

{\setlength{\parindent}{0pt}
\setlength{\parskip}{5pt}
\hangindent=1.27cm
Gray, H. M., Gray, K., \& Wegner, D. M. (2007). Dimensions of mind perception. \textit{Science, 315}(5812), 619. \url{https://doi.org/10.1126/science.1134475}\par
\hangindent=1.27cm
Lasseter, J. (1987). Principles of traditional animation applied to 3D computer animation. \textit{ACM SIGGRAPH Computer Graphics, 21}(4), 35--44. \url{https://doi.org/10.1145/37402.37407}\par
}

\end{document}
